\documentclass[11pt, a4paper]{article}

% ── Codificación y Lenguaje ──────────────────────────────────────────────────
\usepackage[utf8]{inputenc}
\usepackage[T1]{fontenc}
\usepackage[spanish, es-tabla]{babel}

% ── Geometría ───────────────────────────────────────────────────────────────
\usepackage[top=2.5cm, bottom=2.5cm, left=2.5cm, right=2.5cm, headheight=24pt]{geometry}

% ── Matemáticas ─────────────────────────────────────────────────────────────
\usepackage{amsmath, amssymb, amsthm}
\usepackage{mathtools}

% ── Colores y Cajas ─────────────────────────────────────────────────────────
\usepackage{xcolor}
\usepackage[most]{tcolorbox}
\tcbuselibrary{skins, breakable, theorems}

% ── Tablas ──────────────────────────────────────────────────────────────────
\usepackage{booktabs}
\usepackage{array}
\usepackage{multirow}
\usepackage{longtable}

% ── Listas ──────────────────────────────────────────────────────────────────
\usepackage{enumitem}

% ── Encabezado y Pie de Página ──────────────────────────────────────────────
\usepackage{fancyhdr}

% ── Secciones ───────────────────────────────────────────────────────────────
\usepackage{titlesec}

% ── Gráficos y Diagramas ────────────────────────────────────────────────────
\usepackage{tikz}
\usetikzlibrary{shapes.geometric, arrows.meta, positioning, calc}
\usepackage{graphicx}

% ── Columnas múltiples ──────────────────────────────────────────────────────
\usepackage{multicol}

% ── Espaciado ───────────────────────────────────────────────────────────────
\usepackage{setspace}
\setlength{\parskip}{4pt}
\setlength{\parindent}{0pt}

% ── Paleta GOP ──────────────────────────────────────────────────────────────
\definecolor{gopblue}{RGB}{31, 78, 121}
\definecolor{gopcyan}{RGB}{70, 130, 180}
\definecolor{gopgreen}{RGB}{0, 100, 0}
\definecolor{gopgreenlight}{RGB}{230, 255, 230}
\definecolor{goporange}{RGB}{200, 100, 0}
\definecolor{goporangelight}{RGB}{255, 245, 210}
\definecolor{gopred}{RGB}{160, 0, 0}
\definecolor{gopredlight}{RGB}{255, 230, 230}
\definecolor{gopgray}{RGB}{90, 90, 90}
\definecolor{gopgraylight}{RGB}{245, 245, 240}
\definecolor{gopbluedef}{RGB}{0, 70, 140}
\definecolor{gopbluedeflight}{RGB}{220, 235, 255}

% ── Hipervínculos ───────────────────────────────────────────────────────────
\usepackage[colorlinks=true, linkcolor=gopblue, urlcolor=gopblue]{hyperref}

% ── Entornos tcolorbox ───────────────────────────────────────────────────────
\newtcolorbox{definicion}[1]{
  enhanced, breakable,
  colback=gopbluedeflight, colframe=gopbluedef,
  fonttitle=\bfseries\small, title={Definición: #1}, coltitle=white,
  attach boxed title to top left={yshift=-2mm, xshift=6pt},
  boxed title style={colback=gopbluedef, colframe=gopbluedef, sharp corners},
  sharp corners=south, top=4mm, before skip=6pt, after skip=6pt,
}
\newtcolorbox{formulas}[1]{
  enhanced, breakable,
  colback=gopgreenlight, colframe=gopgreen,
  fonttitle=\bfseries\small, title={Fórmulas: #1}, coltitle=white,
  attach boxed title to top left={yshift=-2mm, xshift=6pt},
  boxed title style={colback=gopgreen, colframe=gopgreen, sharp corners},
  sharp corners=south, top=4mm, before skip=6pt, after skip=6pt,
}
\newtcolorbox{tipprueba}[1]{
  enhanced, breakable,
  colback=goporangelight, colframe=goporange,
  fonttitle=\bfseries\small, title={TIP PRUEBA: #1}, coltitle=white,
  attach boxed title to top left={yshift=-2mm, xshift=6pt},
  boxed title style={colback=goporange, colframe=goporange, sharp corners},
  sharp corners=south, top=4mm, before skip=6pt, after skip=6pt,
}
\newtcolorbox{trampa}[1]{
  enhanced, breakable,
  colback=gopredlight, colframe=gopred,
  fonttitle=\bfseries\small, title={TRAMPA/ADVERTENCIA: #1}, coltitle=white,
  attach boxed title to top left={yshift=-2mm, xshift=6pt},
  boxed title style={colback=gopred, colframe=gopred, sharp corners},
  sharp corners=south, top=4mm, before skip=6pt, after skip=6pt,
}
\newtcolorbox{ejercicio}[1]{
  enhanced, breakable,
  colback=gopgraylight, colframe=gopgray,
  fonttitle=\bfseries\small\color{black}, title={Ejercicio: #1},
  attach boxed title to top left={yshift=-2mm, xshift=6pt},
  boxed title style={colback=gopgray!40, colframe=gopgray, sharp corners},
  sharp corners=south, top=4mm, before skip=6pt, after skip=6pt,
}

% ── Encabezado y pie ─────────────────────────────────────────────────────────
\pagestyle{fancy}
\fancyhf{}
\fancyhead[L]{\small\textbf{GOP ICS3213} --- Solución Interrogación 2 (Para la Casa)}
\fancyhead[R]{\small\thepage}
\fancyfoot[L]{\footnotesize Solución Desarrollada --- Transcripción en LaTeX}
\fancyfoot[R]{\small\thepage}
\renewcommand{\headrulewidth}{0.4pt}
\renewcommand{\footrulewidth}{0.4pt}

% ── Estilo de secciones ──────────────────────────────────────────────────────
\titleformat{\section}
  {\color{gopblue}\Large\bfseries}{\color{gopblue}\thesection.}{0.5em}{}
  [\color{gopblue}\titlerule]
\titleformat{\subsection}
  {\color{gopblue}\normalsize\bfseries}{\color{gopblue}\thesubsection.}{0.5em}{}

\begin{document}

% ════════════════════════════════════════════════════════════════════════════
% PORTADA
% ════════════════════════════════════════════════════════════════════════════
\begin{titlepage}
  \centering
  \vspace*{2cm}
  {\Huge\bfseries\color{gopcyan} Solución Oficial Desarrollo\par}
  \vspace{0.4cm}
  {\LARGE\bfseries Interrogación 2 (Para la Casa)\par}
  \vspace{0.3cm}
  {\large Gestión de Operaciones (ICS3213) --- 1er Semestre 2026\par}
  \vspace{1.2cm}
  \rule{\textwidth}{0.4pt}
  \vspace{0.4cm}
  {\small Profesor:\par}
  {\large\bfseries Alejandro Mac Cawley\par}
  \vspace{0.3cm}
  {\small Tutor Auxiliar:\par}
  {\small\textbf{Antigravity AI Tutor}\par}
  \vspace{0.4cm}
  \rule{\textwidth}{0.4pt}
  \vfill
\end{titlepage}

\newpage

% ════════════════════════════════════════════════════════════════════════════
% DESARROLLO DE PREGUNTAS
% ════════════════════════════════════════════════════════════════════════════

\section{Pregunta 1: Planificación Agregada y MRP}

\begin{ejercicio}{Planificación Agregada y MRP --- AgriDrone S.A.}
  \begin{tipprueba}{Estrategia de Nivelación sin Faltantes}
    \textbf{Identificación:} Se solicita determinar una fuerza de trabajo constante ($W$) y un plan de producción mensual constante ($P$) para un horizonte de 3 meses de manera que no existan faltantes al final de cada período. \\
    \textbf{Qué aplicar:} La condición matemática exige que la producción acumulada al final de cada mes cubra la demanda acumulada. Dado que el tamaño de la fuerza laboral y la tasa de producción son constantes, despejamos el valor límite inferior para $P$.
  \end{tipprueba}

  \textbf{Datos de Planificación Agregada:}
  \begin{itemize}
    \item Demanda Octubre ($D_1$): $800$ unidades.
    \item Demanda Noviembre ($D_2$): $1.200$ unidades.
    \item Demanda Diciembre ($D_3$): $1.600$ unidades.
    \item Productividad: $10$ drones/semana por trabajador.
    \item Semanas por mes: $4$ semanas (mes estándar), por ende, la productividad mensual es $40$ drones/mes por trabajador.
    \item Inventario Inicial ($I_0$): $0$ unidades (a fines de septiembre).
    \item Costo de mantener inventario ($H$): $\$50$ por dron al mes (aplicado sobre el inventario final de cada mes).
  \end{itemize}

  \medskip
  \textbf{Desarrollo i. Plan Agregado y Cantidad de Personal:}

  Sea $P$ la producción mensual constante. Se deben cumplir las restricciones de no faltantes acumuladas período a período:
  \begin{align*}
    t = 1 \quad (\text{Oct}): \quad & I_1 = I_0 + P - D_1 = P - 800 \ge 0 \implies P \ge 800 \\
    t = 2 \quad (\text{Nov}): \quad & I_2 = I_1 + P - D_2 = 2P - 2.000 \ge 0 \implies P \ge 1.000 \\
    t = 3 \quad (\text{Dic}): \quad & I_3 = I_2 + P - D_3 = 3P - 3.600 \ge 0 \implies P \ge 1.200
  \end{align*}

  Para satisfacer la demanda del trimestre sin quiebres de stock en ningún período, la producción mensual constante requerida debe ser:
  \[
    P^* = \max(800, 1.000, 1.200) = \mathbf{1.200 \text{ unidades/mes}}
  \]

  La fuerza laboral constante ($W$) necesaria se calcula dividiendo la producción requerida entre la productividad de un trabajador:
  \[
    W = \frac{P^*}{\text{Productividad Mensual}} = \frac{1.200 \text{ unidades/mes}}{40 \text{ unidades/trabajador-mes}} = \mathbf{30 \text{ trabajadores}}
  \]

  Calculamos el balance de inventarios y los costos de almacenamiento al final de cada mes:
  \begin{itemize}
    \item \textbf{Octubre:}
      \[
        I_1 = 0 + 1.200 - 800 = 400 \text{ unidades} \implies \text{Costo} = 400 \times \$50 = \$20.000
      \]
    \item \textbf{Noviembre:}
      \[
        I_2 = 400 + 1.200 - 1.200 = 400 \text{ unidades} \implies \text{Costo} = 400 \times \$50 = \$20.000
      \]
    \item \textbf{Diciembre:}
      \[
        I_3 = 400 + 1.200 - 1.600 = 0 \text{ unidades} \implies \text{Costo} = 0 \times \$50 = \$0
      \]
  \end{itemize}

  El costo total de mantener inventario es:
  \[
    \text{Costo Total Inventario} = \$20.000 + \$20.000 + \$0 = \mathbf{\$40.000}
  \]

  \rule{\textwidth}{0.2pt}

  \textbf{Desarrollo ii. Planificación MRP para MCM (Octubre):}
  \begin{itemize}
    \item La producción agregada de drones en octubre es de $1.200$ unidades.
    \item El **modelo estrella** representa el $50\%$ del total de producción: $1.200 \times 0.5 = 600$ drones estrella al mes.
    \item La producción se distribuye uniformemente en las 4 semanas de octubre: $600 / 4 = 150$ drones estrella por semana en el MPS.
    \item Explosión del BOM: Cada dron requiere exactamente 1 MCM. Los requerimientos brutos ($GR$) de MCM son de $150$ unidades/semana.
    \item Lead Time ($LT$): 2 semanas. Política de loteo: L4L. Inventario inicial de MCM: 50. Recepción programada: 200 en Semana 1.
  \end{itemize}

  \begin{center}
    \small
    \begin{tabular}{lccccc}
      \toprule
      \textbf{Concepto / Semana} & \textbf{Semana 0} & \textbf{Semana 1} & \textbf{Semana 2} & \textbf{Semana 3} & \textbf{Semana 4} \\
      \midrule
      Requerimiento Bruto ($GR$) & & 150 & 150 & 150 & 150 \\
      Recepciones Programadas ($SR$) & & 200 & 0 & 0 & 0 \\
      Stock Disponible Proyectado ($I$) & 50 & 100 & 0 & 0 & 0 \\
      Requerimientos Netos ($NR$) & & 0 & 50 & 150 & 150 \\
      Recepciones Planificadas ($POR$) & & 0 & 50 & 150 & 150 \\
      Lanzamientos Planificados ($PORelease$) & 50 & 150 & 150 & 0 & 0 \\
      \bottomrule
    \end{tabular}
  \end{center}

  \textbf{Detalle del cálculo semana a semana:}
  \begin{itemize}
    \item \textbf{Semana 1:} Inventario disponible final = $50 + 200 - 150 = 100 \ge 0$. No se requiere pedido neto.
    \item \textbf{Semana 2:} Inventario disponible inicial = $100$. Requerimiento bruto = $150$. Requerimiento neto = $150 - 100 = 50$. Con loteo L4L, $POR = 50$. Dado un $LT=2$, el lanzamiento planificado se debió realizar en la **Semana 0** (fines de septiembre).
    \item \textbf{Semana 3:} Inventario disponible inicial = $0$. Requerimiento bruto = $150$. Requerimiento neto = $150$. $POR = 150$. Con $LT=2$, el pedido se debe lanzar ($PORelease$) en la **Semana 1**.
    \item \textbf{Semana 4:} Inventario disponible inicial = $0$. Requerimiento bruto = $150$. Requerimiento neto = $150$. $POR = 150$. Con $LT=2$, el pedido se debe lanzar ($PORelease$) en la **Semana 2**.
  \end{itemize}

  \textbf{Lanzamientos de pedidos requeridos en octubre:}
  \begin{itemize}
    \item En la **Semana 1**, se deben lanzar **150 unidades** de MCM.
    \item En la **Semana 2**, se deben lanzar **150 unidades** de MCM.
  \end{itemize}

  \rule{\textwidth}{0.2pt}

  \textbf{Desarrollo iii. Análisis Crítico EOQ vs. MRP:}

  El modelo **EOQ clásico** asume que la demanda es continua, conocida y constante a lo largo del tiempo. En un sistema MRP, sin embargo, los requerimientos brutos de componentes (demanda dependiente) son de naturaleza **discreta, altamente variable e intermitente (*lumpy demand*)** debido a los desfases del lead time y la explosión en lotes del nivel superior.

  Utilizar EOQ en un MRP genera graves ineficiencias:
  \begin{enumerate}
    \item **Exceso de Inventario:** Al forzar un lote fijo ($Q^*_{\text{EOQ}}$), en períodos de baja demanda se generará stock remanente innecesario que se almacena, incrementando los costos de mantener.
    \item **Retrasos de Producción:** En períodos de alta demanda acumulada en una sola semana ( setup concentrado), el lote fijo puede no ser suficiente para cubrir el requerimiento neto, provocando quiebres de stock.
  \end{enumerate}

  \textbf{Recomendación:} Se aconseja utilizar el algoritmo de **Wagner-Whitin (WW)** o la heurística de **Silver-Meal (SM)**, ya que equilibran de forma dinámica los costos de setup y almacenamiento considerando la demanda discreta y variable de cada período.
\end{ejercicio}

\newpage

\section{Pregunta 2: Administración de Proyectos PERT/CPM}

\begin{ejercicio}{Administración de Proyectos --- Redd SpA}
  \begin{tipprueba}{PERT Estocástico y Evaluación de Riesgo}
    \textbf{Identificación:} Se cuenta con una lista de actividades, duraciones esperadas ($TE$) y desviaciones estándar ($\sigma_i$). Se requiere calcular la ruta crítica, probabilidad de cumplimiento a una fecha límite, e intervalos de confianza. \\
    \textbf{Qué aplicar:} Pasada hacia adelante para obtener los tiempos tempranos ($ES, EF$), pasada hacia atrás para obtener los tiempos tardíos ($LS, LF$), identificar actividades con holgura cero (Ruta Crítica), y aplicar la normal estándar $Z = \frac{X - \mu}{\sigma_p}$ sumando las varianzas de la ruta crítica.
  \end{tipprueba}

  \textbf{Datos de las Actividades (Tiempos en semanas):}
  \begin{center}
    \small
    \begin{tabular}{cccccc}
      \toprule
      Actividad & Predecesor & Tiempo Medio ($TE$) & Desv. Estándar ($\sigma$) & Varianza ($\sigma^2$) \\
      \midrule
      A & - & 3 & 0.667 & 4/9 $\approx$ 0.444 \\
      B & A & 5 & 0.667 & 4/9 $\approx$ 0.444 \\
      C & A & 4 & 1.000 & 1 \\
      D & B & 6 & 1.333 & 16/9 $\approx$ 1.778 \\
      E & B, C & 5 & 0.667 & 4/9 $\approx$ 0.444 \\
      F & C & 7 & 1.000 & 1 \\
      G & D, E & 4 & 0.667 & 4/9 $\approx$ 0.444 \\
      H & F, G & 3 & 0.333 & 1/9 $\approx$ 0.111 \\
      \bottomrule
    \end{tabular}
  \end{center}

  \medskip
  Para determinar la ruta crítica y las holguras, realizamos las dos pasadas del método CPM:
  \begin{itemize}
    \item \textbf{Pasada hacia adelante (ES, EF):} Determina los tiempos de inicio y término más tempranos:
      \begin{align*}
        & ES_A = 0 \implies EF_A = 3 \\
        & ES_B = EF_A = 3 \implies EF_B = 3 + 5 = 8 \\
        & ES_C = EF_A = 3 \implies EF_C = 3 + 4 = 7 \\
        & ES_D = EF_B = 8 \implies EF_D = 8 + 6 = 14 \\
        & ES_E = \max(EF_B, EF_C) = \max(8, 7) = 8 \implies EF_E = 8 + 5 = 13 \\
        & ES_F = EF_C = 7 \implies EF_F = 7 + 7 = 14 \\
        & ES_G = \max(EF_D, EF_E) = \max(14, 13) = 14 \implies EF_G = 14 + 4 = 18 \\
        & ES_H = \max(EF_F, EF_G) = \max(14, 18) = 18 \implies EF_H = 18 + 3 = 21
      \end{align*}
      La duración esperada del proyecto es $E[T] = \mathbf{21 \text{ semanas}}$.

    \item \textbf{Pasada hacia atrás (LS, LF):} Determina los tiempos de inicio y término más tardíos para no retrasar el proyecto (con tiempo final fijado en $LF_H = 21$ semanas):
      \begin{align*}
        & LF_H = 21 \implies LS_H = 21 - 3 = 18 \\
        & LF_G = LS_H = 18 \implies LS_G = 18 - 4 = 14 \\
        & LF_F = LS_H = 18 \implies LS_F = 18 - 7 = 11 \\
        & LF_E = LS_G = 14 \implies LS_E = 14 - 5 = 9 \\
        & LF_D = LS_G = 14 \implies LS_D = 14 - 6 = 8 \\
        & LF_C = \min(LS_E, LS_F) = \min(9, 11) = 9 \implies LS_C = 9 - 4 = 5 \\
        & LF_B = \min(LS_D, LS_E) = \min(8, 9) = 8 \implies LS_B = 8 - 5 = 3 \\
        & LF_A = \min(LS_B, LS_C) = \min(3, 5) = 3 \implies LS_A = 3 - 3 = 0
      \end{align*}
  \end{itemize}

  \textbf{Tabla Consolidada de Tiempos y Holguras (CPM):}
  Se presenta a continuación el resumen estructurado de las pasadas hacia adelante y hacia atrás. La holgura ($H$) se calcula como $H_i = LS_i - ES_i = LF_i - EF_i$.
  \begin{center}
    \small
    \begin{tabular}{ccccccccc}
      \toprule
      Actividad & Predecesor & Duración ($TE$) & $ES$ & $EF$ & $LS$ & $LF$ & Holgura ($H$) & Crítica \\
      \midrule
      A & - & 3 & 0 & 3 & 0 & 3 & 0 & Sí \\
      B & A & 5 & 3 & 8 & 3 & 8 & 0 & Sí \\
      C & A & 4 & 3 & 7 & 5 & 9 & 2 & No \\
      D & B & 6 & 8 & 14 & 8 & 14 & 0 & Sí \\
      E & B, C & 5 & 8 & 13 & 9 & 14 & 1 & No \\
      F & C & 7 & 7 & 14 & 11 & 18 & 4 & No \\
      G & D, E & 4 & 14 & 18 & 14 & 18 & 0 & Sí \\
      H & F, G & 3 & 18 & 21 & 18 & 21 & 0 & Sí \\
      \bottomrule
    \end{tabular}
  \end{center}

  Por lo tanto, la **Ruta Crítica** está conformada por las actividades **A - B - D - G - H**.
  \begin{itemize}
    \item Tiempos del proyecto: $E[T] = \mathbf{21 \text{ semanas}}$.
    \item Varianza del proyecto ($\text{Var}[T]$):
      \[
        \text{Var}[T] = \sigma_A^2 + \sigma_B^2 + \sigma_D^2 + \sigma_G^2 + \sigma_H^2 = \frac{4}{9} + \frac{4}{9} + \frac{16}{9} + \frac{4}{9} + \frac{1}{9} = \frac{29}{9} \approx \mathbf{3.222 \text{ semanas}^2}
      \]
    \item Desviación estándar del proyecto ($\sigma_p$):
      \[
        \sigma_p = \sqrt{3.222} \approx \mathbf{1.795 \text{ semanas}}
      \]
  \end{itemize}

  \rule{\textwidth}{0.2pt}

  \textbf{Desarrollo ii. Probabilidad e Intervalo de Confianza:}
  \begin{itemize}
    \item Probabilidad de terminar en o antes de 23 semanas ($P(T \le 23)$):
      \[
        Z = \frac{23 - E[T]}{\sigma_p} = \frac{23 - 21}{1.795} = 1.114
      \]
      Buscando en la tabla de la normal estándar para $Z = 1.11$, obtenemos:
      \[
        \Phi(1.11) \approx 0.8665 \implies P(T \le 23) \approx \mathbf{86.65\%}
      \]
      *(Utilizando interpolación lineal para $Z = 1.114$: $\Phi(1.114) \approx 0.8665 + 0.4 \times (0.8686 - 0.8665) = \mathbf{86.73\%}$)*

    \item Intervalo de confianza al 95\% (dos colas):
      Para el 95\% de confianza, la cota crítica es $Z_{\alpha/2} = 1.96$:
      \[
        IC = E[T] \pm 1.96 \cdot \sigma_p = 21 \pm 1.96 \cdot (1.795) = 21 \pm 3.518 \implies \mathbf{[17.482, 24.518] \text{ semanas}}
      \]
  \end{itemize}

  \rule{\textwidth}{0.2pt}

  \textbf{Desarrollo iii. Evaluación de Contrato (Bono y Penalización Fija):}
  \begin{itemize}
    \item Bono: $\$1.5$ millones si $T \le 21$.
      \[
        Z_{\text{bono}} = \frac{21 - 21}{\sigma_p} = 0 \implies P(T \le 21) = 0.50
      \]
    \item Penalidad: $\$4.0$ millones si $T > 23$.
      \[
        P(T > 23) = 1 - P(T \le 23) = 1 - 0.8673 = 0.1327
      \]
    \item Valor Esperado del Contrato ($VE$):
      \[
        VE = 1.5 \cdot P(T \le 21) - 4.0 \cdot P(T > 23) = 1.5(0.50) - 4.0(0.1327) = 0.75 - 0.5308 = \mathbf{+0.2192 \text{ millones}}
      \]
  \end{itemize}
  **Conclusión:** Dado que el valor esperado de la ganancia neta del contrato es positivo ($+\$219.200$), **se debe aceptar el contrato**.

  \rule{\textwidth}{0.2pt}

  \textbf{Desarrollo iv. Modelo Matemático para Bono/Penalidad Lineal:}

  Si se penaliza con $\$4.0$ millones por cada semana de atraso ($T > 23$) y se bonifica con $\$1.5$ millones por semana de adelanto ($T \le 21$), el resultado financiero del contrato $f(T)$ es una función continua por tramos de la variable aleatoria de finalización $T \sim N(21, 3.222)$:
  \[
    f(T) = \begin{cases}
      1.5 \cdot (21 - T) & \text{si } T \le 21 \\
      0 & \text{si } 21 < T \le 23 \\
      -4.0 \cdot (T - 23) & \text{si } T > 23
    \end{cases}
  \]

  La condición matemática para aceptar el contrato es que su valor esperado sea positivo ($E[f(T)] > 0$). Expresado de forma analítica en función de la densidad de probabilidad $\phi_T(t)$:
  \[
    \int_{-\infty}^{21} 1.5(21 - t) \frac{1}{\sqrt{2\pi}\sigma_p} e^{-\frac{(t-21)^2}{2\sigma_p^2}} dt - \int_{23}^{\infty} 4.0(t - 23) \frac{1}{\sqrt{2\pi}\sigma_p} e^{-\frac{(t-21)^2}{2\sigma_p^2}} dt > 0
  \]
\end{ejercicio}

\newpage

\section{Pregunta 3: Gestión de Capacidad y Variabilidad}

\begin{ejercicio}{Gestión de Capacidad y Variabilidad --- Registro Civil}
  \begin{tipprueba}{Aproximación de Kingman VUT}
    \textbf{Identificación:} Se dispone de tiempos de llegada y atención con sus respectivos coeficientes de variación de arribos ($C_a$) y servicios ($C_s$). Se busca calcular el tiempo de espera en cola ($W_q$). \\
    \textbf{Qué aplicar:} La aproximación de Kingman para colas generales $G/G/1$: $W_q = \left(\frac{C_a^2 + C_s^2}{2}\right) \left(\frac{\rho}{1-\rho}\right) \left(\frac{1}{\mu}\right)$, y luego Little $L_q = \lambda W_q$.
  \end{tipprueba}

  \textbf{Datos Iniciales:}
  \begin{itemize}
    \item Tasa de llegadas ($\lambda$): $20$ personas/hora $\implies \lambda = 1/3$ personas/minuto.
    \item Tiempo promedio de servicio ($1/\mu$): $2.4$ minutos $\implies \mu = 25$ personas/hora.
    \item Coeficiente de variación de arribos ($C_a$): $0.5$.
    \item Coeficiente de variación del servicio ($C_s$): $1.5$.
  \end{itemize}

  \medskip
  \textbf{Desarrollo a) Utilización y Tiempos de Espera:}

  El nivel de utilización de la oficina ($\rho$) se define como:
  \[
    \rho = \frac{\lambda}{\mu} = \frac{20 \text{ personas/hora}}{25 \text{ personas/hora}} = \mathbf{0.80 \quad (80\%)}
  \]

  El tiempo medio de espera en cola ($W_q$) se aproxima mediante la ecuación de Kingman:
  \[
    W_q = \left( \frac{C_a^2 + C_s^2}{2} \right) \left( \frac{\rho}{1 - \rho} \right) \left( \frac{1}{\mu} \right)
  \]
  Reemplazando los valores del enunciado:
  \[
    W_q = \left( \frac{0.5^2 + 1.5^2}{2} \right) \left( \frac{0.80}{1 - 0.80} \right) (2.4 \text{ minutos})
  \]
  \[
    W_q = \left( \frac{0.25 + 2.25}{2} \right) \left( \frac{0.80}{0.20} \right) (2.4) = (1.25) \cdot (4) \cdot (2.4) = \mathbf{12 \text{ minutos}}
  \]
  Un cliente espera en promedio **12 minutos** en la cola.

  El tamaño promedio de la sala de espera (cantidad de clientes en la cola, $L_q$) se obtiene por la Ley de Little:
  \[
    L_q = \lambda \cdot W_q = 20 \text{ personas/hora} \times \frac{12 \text{ minutos}}{60 \text{ min/hora}} = 20 \times 0.2 = \mathbf{4 \text{ personas}}
  \]
  El tamaño promedio de asientos necesarios en la sala de espera es de **4 asientos**.

  \rule{\textwidth}{0.2pt}

  \textbf{Desarrollo b) Agenda Digital en Bloques de 1 Hora:}

  Si se implementa una agenda digital en bloques regulares de 1 hora, se elimina la variabilidad aleatoria de los tiempos entre arribos de los clientes, haciendo que las llegadas ocurran en intervalos fijos predeterminados. Esto implica que el coeficiente de variación de los arribos es nulo: $C_a = 0$ (arribos determinísticos).

  Buscamos la tasa máxima de agendamientos $\lambda$ (y por ende la utilización $\rho$) para garantizar que el tiempo de espera no exceda los 6 minutos ($W_q \le 6$):
  \[
    W_q = \left( \frac{0^2 + 1.5^2}{2} \right) \left( \frac{\rho}{1 - \rho} \right) \left( 2.4 \right) \le 6
  \]
  \[
    1.125 \cdot 2.4 \cdot \left( \frac{\rho}{1 - \rho} \right) \le 6 \implies 2.7 \left( \frac{\rho}{1 - \rho} \right) \le 6
  \]
  Despejamos la relación de utilización:
  \[
    \frac{\rho}{1 - \rho} \le \frac{6}{2.7} = \frac{20}{9} \implies 9\rho \le 20(1-\rho) \implies 29\rho \le 20 \implies \rho \le \frac{20}{29} \approx 0.6897
  \]

  Utilizando la definición de utilización, calculamos la tasa máxima de llegada $\lambda$:
  \[
    \lambda = \rho \cdot \mu \le 0.6897 \cdot 25 = 17.24 \text{ clientes/hora}
  \]
  Por lo tanto, la cantidad máxima de clientes que se puede agendar por hora para asegurar que la espera sea menor a 6 minutos es de **17 clientes**.

  \rule{\textwidth}{0.2pt}

  \textbf{Desarrollo c) Evaluación de Opciones de Mejora (Cs vs. Tasa de Llegada):}

  Sabemos que en la situación base: $W_{q,\text{inicial}} = 12 \text{ minutos}$.
  El beneficio neto de cada opción se define como el beneficio por reducción de espera total ($\Delta W_q \times \$1.000$) menos el costo del proyecto.

  \begin{itemize}
    \item \textbf{Opción i: Capacitación al Personal ($C_s \to 1.0$)}
      \begin{itemize}
        \item Parámetros: $\lambda = 20, \mu = 25 \implies \rho = 0.80$. $C_a = 0.5$, nuevo $C_s = 1.0$.
        \item Nuevo tiempo de espera en cola:
          \[
            W_{q,\text{Cap}} = \left( \frac{0.5^2 + 1.0^2}{2} \right) \left( \frac{0.80}{1 - 0.80} \right) (2.4) = (0.625) \cdot (4) \cdot (2.4) = \mathbf{6 \text{ minutos}}
          \]
        \item Reducción de tiempo ($\Delta W_q$): $12 - 6 = 6 \text{ minutos}$.
        \item Beneficio bruto: $6 \text{ min} \times \$1.000/\text{min} = \$6.000$.
        \item Costo: $\$4.000$.
        \item Beneficio neto: $\$6.000 - \$4.000 = \mathbf{+\$2.000}$.
      \end{itemize}

    \item \textbf{Opción ii: Digitalización de Servicios ($\lambda \to 15$)}
      \begin{itemize}
        \item Parámetros: nuevo $\lambda = 15, \mu = 25 \implies \rho = 0.60$. $C_a = 0.5$, $C_s = 1.5$.
        \item Nuevo tiempo de espera en cola:
          \[
            W_{q,\text{Dig}} = \left( \frac{0.5^2 + 1.5^2}{2} \right) \left( \frac{0.60}{1 - 0.60} \right) (2.4) = (1.25) \cdot (1.5) \cdot (2.4) = \mathbf{4.5 \text{ minutos}}
          \]
        \item Reducción de tiempo ($\Delta W_q$): $12 - 4.5 = 7.5 \text{ minutos}$.
        \item Beneficio bruto: $7.5 \text{ min} \times \$1.000/\text{min} = \$7.500$.
        \item Costo: $\$6.000$.
        \item Beneficio neto: $\$7.500 - \$6.000 = \mathbf{+\$1.500}$.
      \end{itemize}
  \end{itemize}

  **Conclusión:** Se prefiere la **Opción i** (Capacitación al personal), ya que reporta un beneficio neto mayor de **$+\$2.000$** en comparación con los $+\$1.500$ de la Opción ii.
\end{ejercicio}

\end{document}
